Los resultados indican una brecha estructural entre sedes en condiciones de estudio. Esta variabilidad puede explicarse por diferencias en diseño arquitectónico, antigüedad de infraestructura, políticas de operación local, densidad de usuarios y capacidades de mantenimiento.

Desde la perspectiva de infraestructura pública urbana, los hallazgos sugieren que la calidad del servicio bibliotecario no depende solo de conectividad, sino de la interacción entre ambiente físico (ruido e iluminación), gobernanza del espacio y gestión por demanda. En contextos de alta ocupación, la calidad percibida por usuarios puede deteriorarse aun cuando la red WiFi sea aceptable.

Para política pública, esto implica migrar de un enfoque de cumplimiento mínimo a uno de desempeño integral por experiencia de usuario. En particular, el criterio acústico emerge como cuello de botella transversal, con impacto potencial sobre concentración, permanencia y rendimiento académico.
